\section{A longitudinal case: the 2019 unified framework}

In 2019, Dauparas et al. proposed a \emph{Unified Framework for Modeling Multivariate
Distributions in Biological Sequences} \cite{dauparas2019}.  The paper placed PSSMs, Potts or
Markov random fields, multivariate Gaussian models, and autoencoders inside one reconstruction
diagram.  Read retrospectively, its importance lies in two achievements.  It exposed exact
equivalences that had been hidden by disciplinary notation, and it made those equivalences
operational enough that a regularization choice could be transferred from one model family to
another.  It also shows where semantic typing begins after an algebraic unification succeeds.

This distinction marks a genuine progression that builds on the original result.  The
2019 question was: how many familiar sequence models can be implemented as nearby versions of one
input--output reconstruction system?  The question of this book is stricter: after that
implementation has been unified, which mathematical object does each system identify, and what
must remain different when probability, energy, computation, and experiment are translated?

\subsection{The common reconstruction diagram}

Let an MSA contain $N$ sequences of length $L$ over an alphabet of size $q$.  Flatten the one-hot
tensor to $X\in\{0,1\}^{N\times d}$ with $d=Lq$.  A linear reconstruction system has preactivation
\begin{equation}
H=B+XW,
\label{eq:book-unified-linear-map}
\end{equation}
where $B$ supplies one-site fields and $W$ carries cross-coordinate dependence.  The models in the
2019 framework are obtained by changing four choices: the admissible blocks of $W$, the output
loss, the treatment of self-connections, and the regularizer.  Adding hidden layers replaces the
single linear map by a nonlinear encoder--decoder while preserving the reconstruction view.

For a PSSM, $W=0$ and a sitewise categorical softmax learns only $B_i(a)$.  For a Potts
pseudolikelihood model, the logits at target position $i$ are
\begin{equation}
H_i(a;x_{-i})
=B_i(a)+\sum_{j\ne i}\sum_b x_j(b)W_{j\to i}(b,a),
\label{eq:book-unified-potts-logit}
\end{equation}
and categorical cross-entropy is applied over the residue alphabet at $i$.  Because one-hot targets
satisfy $\sum_a x_i(a)=1$, the loss is exactly the negative sitewise Potts log
pseudolikelihood:
\begin{equation}
\mathcal L_{\mathrm{CCE}}
=-\frac1N\sum_{n,i}
\log
\frac{\exp H_i(x_i^{(n)};x_{-i}^{(n)})}
{\sum_a\exp H_i(a;x_{-i}^{(n)})}
=\mathcal L_{\mathrm{PL}}.
\label{eq:book-unified-cce-pl}
\end{equation}
The equality is an exact identity of objectives after the categorical sample space and the
targetwise normalization have been matched.

The Gaussian branch begins from a different law.  Center the flattened observations to
$\widetilde X$ and consider a regularized self-reconstruction objective of the form
\begin{equation}
\mathcal L_{\mathrm{MSE}}(W)
=\frac1N\|\widetilde X-\widetilde XW\|_F^2
+\lambda\|W\|_F^2-2\gamma\operatorname{tr}(W).
\label{eq:book-unified-mse}
\end{equation}
Writing $C=N^{-1}\widetilde X^{\mathsf T}\widetilde X$, its stationary solution is
\begin{equation}
W=(\gamma-\lambda)(C+\lambda I)^{-1}+I.
\label{eq:book-unified-mse-solution}
\end{equation}
After the identity map is separated, the learned cross-coordinate map is therefore an affine
transform of a regularized precision matrix.  With the corresponding choice of $\gamma$, the MSE
system and Gaussian likelihood share the same global optimum, even though one is written as a
neural self-map and the other as covariance inversion.

The resulting family resemblance can be summarized without erasing the differences:
\begin{center}
\small
\begin{tabular}{@{}>{\raggedright\arraybackslash}p{0.14\linewidth}
                >{\raggedright\arraybackslash}p{0.19\linewidth}
                >{\raggedright\arraybackslash}p{0.16\linewidth}
                >{\raggedright\arraybackslash}p{0.19\linewidth}
                >{\raggedright\arraybackslash}p{0.19\linewidth}@{}}
\toprule
Model & Reconstruction map & Objective & Native probability object & First semantic remainder \\
\midrule
PSSM & bias field only & sitewise CCE & independent categorical sites & no represented pair edge \\
Potts/MRF & linear cross-position blocks & categorical CCE = pseudolikelihood & compatible categorical conditionals, when parameters are tied into one law & raw block retains gauge freedom; contact is a separate endpoint \\
Gaussian/MG & centered linear self-map & regularized MSE & continuous Gaussian law through precision & one-hot simplex and categorical contrasts are approximated \\
Autoencoder & factorized or nonlinear reconstruction & chosen output loss & generally a conditional predictor unless a joint law is supplied & hidden paths can induce higher-order final responses \\
\bottomrule
\end{tabular}
\end{center}

The diagram unified implementation at exactly the moment when neural-network software made that
unification useful.  Pseudolikelihood Potts fitting could be expressed as ordinary categorical
reconstruction; Gaussian inverse-covariance estimation could be expressed through an MSE layer;
and either branch could acquire hidden layers without changing the outer input--output task.

\subsection{The empirical test made the unification scientific}

The 2019 work went beyond redrawing four models with the same boxes.  Its contact-prediction
experiment asked whether a choice that worked in one formulation could improve another.  The MRF
literature commonly used $L_2$ regularization, whereas Gaussian DCA was often expressed through a
pseudocount applied to the covariance.  Transferring the MRF-style shrinkage logic to the MSE
formulation improved contact prediction over GaussDCA.  Removing the diagonal self-map gave a
further small improvement and connected the resulting score to partial correlation.  The
categorical cross-entropy model still performed better, but the gap narrowed when the
regularization comparison was made more like-for-like.

This is the strongest form of useful mathematical unification.  An equivalence should permit a
controlled transfer of an estimator, regularizer, diagnostic, or computational method.  Here it
did.  The result showed that part of the apparent difference between model families came from
training conventions and could be reduced by matching those conventions across graphical forms.
It also showed that the remaining difference mattered: categorical cross-entropy respected the
amino-acid sample space more faithfully than squared error on a continuous relaxation.

The contact endpoint introduced one more irreversible step.  Both fitted matrices were reduced to
scalar position scores, corrected by average-product correction, and ranked against structural
contacts.  Success at that endpoint established predictive utility for contact ranking.  Gaussian
precision entries, Potts tables, and neural weights retained their native meanings; the common
benchmark compared their scalar shadows after model-specific estimation and shared postprocessing.

\subsection{What was unified, and what remained plural}

The phrase ``the difference between the models comes down to the loss function'' is exact at the
level of the single-layer reconstruction diagram under the paper's stated constraints.  It becomes
too strong if it is transported unchanged to probability or physics.  Four remainders survive.

First, the sample spaces differ.  Potts pseudolikelihood normalizes over residue identities at one
site while holding a categorical background fixed.  Gaussian MSE treats centered coordinates as
continuous.  The shared outer matrix multiplication therefore sits above two distinct probability
geometries.

Second, a collection of categorical reconstructions yields one joint energy when parameter tying,
gauge handling, and compatibility of the
targetwise conditionals.  A generic linear CCE network may have directed blocks
$W_{j\to i}$ and $W_{i\to j}$ that fail the reciprocity criterion of
Eq.~\ref{eq:reciprocity}.  Its loss remains well defined, while the joint-Hamiltonian reading stops
at the failed compatibility condition.

Third, the matrix $W$ carries several kinds of structure at once.  Its diagonal may encode trivial
self-reconstruction.  Its row and column directions may contain one-body gauge terms.  A block norm
for contact prediction removes sign and categorical orientation.  The 2019 framework made the
shared matrix visible; the reference-centered operator construction of Chapter~3 identifies which
part of each categorical block is irreducibly bivariate.

Fourth, adding hidden layers changes more than expressive power.  Nonlinear composition can turn
pairwise first-layer messages into higher-order output responses.  The weight connecting two
visible coordinates, a hidden computational route, and the audited substitution response of the
complete autoencoder are therefore different edge candidates.  The reconstruction diagram leaves
their selection to a typed output audit.

These remainders can be stated as a progression of audit maps:
\begin{equation}
\begin{aligned}
\text{2019:}\quad
M&\longmapsto
(\text{reconstruction map},\text{ loss},\text{ regularizer}),\\
\text{this book:}\quad
M&\longmapsto
(\text{fiber},\text{ edge operator},\text{ context},\text{ order},
\text{ normalization},\text{ estimator},\text{ validation}).
\end{aligned}
\label{eq:book-two-unifications}
\end{equation}
The second map extends the first from the point where it succeeds.  Once common code
and common algebra have removed superficial differences, the remaining differences become sharp
enough to type.

\subsection{From horizontal equivalence to a longitudinal program}

The 2019 framework is a horizontal unification across contemporary model families.  Its later
consequences are longitudinal.  PSSM bias fields became the simplest node-state baseline.  Potts
CCE made masked categorical prediction continuous with pseudolikelihood.  Autoencoder hidden
layers opened the route from static family models to amortized nonlinear predictors.  Modern pLMs
greatly enlarged that route, but in doing so separated the training objective, internal
communication graph, and query-specific output response more strongly than the one-layer picture
had to confront.

The scientific question consequently moved.  It was first useful to ask whether PSSM, MRF,
Gaussian, and autoencoder fitting could share one computational grammar.  It is now necessary to
ask which typed relations survive depth, pretraining, context dependence, and geometric
supervision.  The operator graph is the next unification in that sequence: a coordinate system
that records the exact correspondence and the exact remainder together.

This case also supplies a standard for the longitudinal analyses that follow.  A historical
sequence of models becomes theoretically informative only when each transition identifies four
things: the state that was added, the information that was discarded, the objective that licensed
the new state, and the scientific claim that still remained unavailable.  The GCN lineage below
applies the standard to graph computation; Chapter~8 then applies it to the three generations of
AlphaFold.
